\chapter{Operator panels}
\label{ch:rviz}

\pkg{dharma_rviz} provides five RViz panels (C++, Qt) to operate and tune a robot, and the layout
\code{rviz/robot.rviz} with one of each next to the map, the planner path and costmap and the
markers of navigation, routes, avoidance, perception and tags. The panels use relative names; the
robot's own topics are remapped to its namespace:
\begin{lstlisting}
rviz2 -d $(ros2 pkg prefix dharma_rviz)/share/dharma_rviz/rviz/robot.rviz \
  --ros-args -r battery:=/robot/battery -r cmd_vel:=/robot/cmd_vel \
  -r motor_temperature:=/robot/motor_temperature
\end{lstlisting}

\section{Panels}

\begin{small}
\begin{longtable}{@{}>{\raggedright}p{0.17\textwidth}>{\raggedright}p{0.39\textwidth}>{\raggedright\arraybackslash}p{0.37\textwidth}@{}}
  \toprule
  Panel & Calls & Shows \\
  \midrule
  \endhead
  \bottomrule
  \endfoot
  Mission (\code{ControlPanel}) & \topic{localization/start}, \topic{route/start},
    \topic{route/stop}, \topic{route/mark_stop_point}, \topic{navigation/load_path},
    \topic{navigation/start}, \topic{stop}, \topic{pause}, \topic{resume},
    \topic{planner/go_to}, \topic{planner/cancel} & \topic{localization/status},
    \topic{navigation/status} (state, progress, pause reasons), \topic{planner/status},
    \topic{route/saved}, the e-stop \\
  Safety (\code{SafetyPanel}) & \topic{control/estop/engage} and \topic{control/estop_request},
    \topic{control/estop/release}, \topic{control/set_mode} & \topic{control/status}: mode,
    remote source, e-stop reasons, armed, protective field, clearance, speed scale \\
  Status (\code{StatusPanel}) & -- & \topic{battery}, \topic{motor_temperature},
    \topic{diagnostics} (worst of all), \topic{health/status}, \topic{odometry/status},
    \topic{cmd_vel} \\
  Tuning (\code{TuningPanel}) & the parameter services of any node & its parameters, ranges and
    read-only flags from the node's descriptors \\
  Tags (\code{TagPanel}) & \topic{localization/record_tags}, \topic{localization/start} &
    \topic{localization/status} \\
\end{longtable}
\end{small}

\section{Behaviour}

\begin{itemize}
  \item \textbf{Never blocking.} Each panel instance has its own node
    (\code{dharma_rviz_<panel>_<pid>_<n>}, review \#10), spun from the Qt event loop. Service
    calls are asynchronous with a client per request and fail with a message after a timeout
    (5~s, the e-stop 2~s) instead of freezing RViz or leaving buttons disabled (\#15, \#16).
  \item \textbf{E-stop first} (\#11). The E-STOP button is never disabled and sends both the
    Trigger service and the \topic{control/estop_request} topic on every press, so it works even
    when one path is down. Start, Resume and Go of the Mission panel are disabled while
    \topic{control/status} reports the e-stop (navigation refuses them anyway); Stop is always
    available. Release is a separate button; the mux still requires the re-arm
    (Section~\ref{sec:safety-case}).
  \item \textbf{Stale values} (\#16). Values that stop arriving for 3~s are greyed out with their
    age; the diagnostics line shows the worst status among those received within that time.
  \item \textbf{Tuning} (\#9, \#15). Edits wait for \emph{Apply} (\emph{Apply on change} is off by
    default); spin boxes take the mouse wheel only when focused; ranges come from
    \code{describe_parameters} and read-only parameters (\code{live=False}) are disabled; choices
    listed in a description (\code{(a | b)}) become a combo box; string arrays are edited one
    item per line; \param{controller} and \param{policy} select which implementation section is
    shown; a second Load replaces the first and its late answers are dropped. \emph{Save YAML}
    writes every shown value at full precision.
  \item \textbf{Goals.} RViz's \emph{2D Goal Pose} tool only fills in the Go To pose of the
    Mission panel; Go sends it to \topic{planner/go_to}.
\end{itemize}

The tests are gtests: the pure helpers (\code{src/common}), the panels in-process with an
offscreen Qt and a fake server node, and every panel loaded through pluginlib
(Appendix~\ref{app:trace}).
